% from ega2-1-fr.tex
\oldpage[II]{17}
ce qui revient au même (\hyperref[0.4.4.3-fr]{0, 4.4.3}) aux
$(\mathcal{O}_S|U)$-homomorphismes
$j^*(\mathcal{E})=\mathcal{E}|U\to\mathcal{O}_S|U$. En outre, il est
immédiat qu'à la restriction à un ouvert $U'\subset U$ d'un $S$-morphisme
$U\to\mathbf{V}(\mathcal{E})$ correspond la restriction à $U'$ de
l'homomorphisme correspondant
$\mathcal{E}|U\to\mathcal{O}_S|U$. On en conclut que \emph{le faisceau des
germes de $S$-sections} de $\mathbf{V}(\mathcal{E})$ s'identifie
canoniquement au \emph{dual $\check{\mathcal{E}}$} de $\mathcal{E}$.

En particulier, si on prend $X=U=S$, à l'homomorphisme \emph{nul}
$\mathcal{E}\to\mathcal{O}_S$ correspond une $S$-section canonique de
$\mathbf{V}(\mathcal{E})$, dite \emph{$S$-section nulle}
(cf. (\hyperref[II.8.3.3-fr]{8.3.3})).
\end{env}

\begin{env}[1.7.10]
\label{II.1.7.10-fr}
Prenons maintenant pour $X$ le spectre $\{\xi\}$ d'un corps $K$~; le
morphisme structural $f:X\to S$ correspond donc à un monomorphisme
$k(s)\to K$, où $s=f(\xi)$
(\hyperref[I.2.4.6-fr]{I, 2.4.6})~; les $S$-morphismes
$\{\xi\}\to\mathbf{V}(\mathcal{E})$ ne sont donc autres que les
\emph{points géométriques de $\mathbf{V}(\mathcal{E})$ à valeurs dans
l'extension $K$ de $k(s)$} (\hyperref[I.3.4.5-fr]{I, 3.4.5}), points qui
sont localisés aux points de $p^{-1}(s)$. L'ensemble de ces points, qu'on
peut encore appeler \emph{la fibre géométrique rationnelle sur $K$ de
$\mathbf{V}(\mathcal{E})$ au-dessus du point $s$}, s'identifie d'après
(\hyperref[II.1.7.8-fr]{1.7.8}) à l'ensemble des homomorphismes de
$\mathcal{O}_S$-Modules
$\mathcal{E}\to f_*(\mathcal{O}_X)$, ou, ce qui revient au même
(\hyperref[0.4.4.3-fr]{0, 4.4.3}) à l'ensemble des homomorphismes de
$\mathcal{O}_X$-Modules
$f^*(\mathcal{E})\to\mathcal{O}_X=K$. Mais on a par définition
(\hyperref[0.4.3.1-fr]{0, 4.3.1})
$f^*(\mathcal{E})=\mathcal{E}_s\otimes_{\mathcal{O}_s}K
=\mathcal{E}^s\otimes_{k(s)}K$, en posant
$\mathcal{E}^s=\mathcal{E}_s/\mathfrak{m}_s\mathcal{E}_s$~; la fibre
géométrique de $\mathbf{V}(\mathcal{E})$ rationnelle sur $K$ au-dessus de
$s$ s'identifie donc au \emph{dual} du \emph{$K$-espace vectoriel
$\mathcal{E}^s\otimes_{k(s)}K$}~; si $\mathcal{E}^s$ ou $K$ est de
dimension finie sur $k(s)$, ce dual s'identifie aussi à
$(\mathcal{E}^s)^\vee\otimes_{k(s)}K$, en désignant par
$(\mathcal{E}^s)^\vee$ le dual du $k(s)$-espace vectoriel
$\mathcal{E}^s$.
\end{env}

\begin{proposition}[1.7.11]
\label{II.1.7.11-fr}
\begin{enumerate}
  \item[\rm(i)] $\mathbf{V}(\mathcal{E})$ est un foncteur contravariant en
    $\mathcal{E}$ de la catégorie des $\mathcal{O}_S$-Modules
    quasi-cohérents dans la catégorie des $S$-schémas affines.
  \item[\rm(ii)] Si $\mathcal{E}$ est un $\mathcal{O}_S$-Module de type
    fini, $\mathbf{V}(\mathcal{E})$ est de type fini sur $S$.
  \item[\rm(iii)] Si $\mathcal{E}$ et $\mathcal{F}$ sont deux
    $\mathcal{O}_S$-Modules quasi-cohérents,
    $\mathbf{V}(\mathcal{E}\oplus\mathcal{F})$ s'identifie canoniquement à
    $\mathbf{V}(\mathcal{E})\times_S\mathbf{V}(\mathcal{F})$.
  \item[\rm(iv)] Soit $g:S'\to S$ un morphisme~; pour tout
    $\mathcal{O}_S$-Module quasi-cohérent $\mathcal{E}$,
    $\mathbf{V}(g^*(\mathcal{E}))$ s'identifie canoniquement à
    $\mathbf{V}(\mathcal{E})_{(S')}
    =\mathbf{V}(\mathcal{E})\times_S S'$.
  \item[\rm(v)] À un homomorphisme surjectif
    $\mathcal{E}\to\mathcal{F}$ de $\mathcal{O}_S$-Modules quasi-cohérents
    correspond une immersion fermée
    $\mathbf{V}(\mathcal{F})\to\mathbf{V}(\mathcal{E})$.
\end{enumerate}
\end{proposition}

\begin{proof}
(i) est conséquence immédiate de
(\hyperref[II.1.2.7-fr]{1.2.7}), compte tenu de ce que tout homomorphisme de
$\mathcal{O}_S$-Modules $\mathcal{E}\to\mathcal{F}$ définit canoniquement un
homomorphisme de $\mathcal{O}_S$-Algèbres
$\mathbf{S}(\mathcal{E})\to\mathbf{S}(\mathcal{F})$. (ii) découle
immédiatement des définitions (\hyperref[I.6.3.1-fr]{I, 6.3.1}) et du fait
que si $E$ est un $A$-module de type fini, $\mathbf{S}(E)$ est une
$A$-algèbre de type fini. Pour démontrer (iii), il suffit de partir de
l'isomorphisme canonique
$\mathbf{S}(\mathcal{E}\oplus\mathcal{F})
\simeq\mathbf{S}(\mathcal{E})
\otimes_{\mathcal{O}_S}\mathbf{S}(\mathcal{F})$
(\hyperref[II.1.7.4-fr]{1.7.4}) et d'appliquer
(\hyperref[II.1.4.6-fr]{1.4.6}). De même, pour démontrer (iv), il suffit de
partir de l'isomorphisme canonique
$\mathbf{S}(g^*(\mathcal{E}))
\simeq g^*(\mathbf{S}(\mathcal{E}))$
(\hyperref[II.1.7.5-fr]{1.7.5}) et d'appliquer
(\hyperref[II.1.5.2-fr]{1.5.2}). Enfin, pour établir (v), il suffit de
remarquer que si l'homomorphisme $\mathcal{E}\to\mathcal{F}$ est surjectif,
il en est de même de l'homomorphisme correspondant
$\mathbf{S}(\mathcal{E})\to\mathbf{S}(\mathcal{F})$ de
$\mathcal{O}_S$-Algèbres, et la conclusion résulte de
(\hyperref[II.1.4.10-fr]{1.4.10}).
\end{proof}

\begin{env}[1.7.12]
\label{II.1.7.12-fr}
Prenons en particulier $\mathcal{E}=\mathcal{O}_S$~; le préschéma
$\mathbf{V}(\mathcal{O}_S)$ est le $S$-schéma affine, spectre de la
$\mathcal{O}_S$-Algèbre $\mathbf{S}(\mathcal{O}_S)$ qui s'identifie à la
$\mathcal{O}_S$-Algèbre
$\mathcal{O}_S[T]=\mathcal{O}_S\otimes_{\mathbf{Z}}\mathbf{Z}[T]$
